\VerdierPhysicalPageAnchor{0077}%
\VerdierSemanticAnchor{ch1-2-4-3-simple-complex-morphisms}
Whence, on passing to the simple complexes, two morphisms of functors:

\VerdierSemanticAnchor{ch1-2-4-3-eq2}
\begin{equation}
  \begin{aligned}
    \operatorname{comp}_{\Pi}h\colon
      \operatorname{comp}_{\Pi}F
      &\longrightarrow \operatorname{comp}_{\Pi}F',\\
    \operatorname{comp}_{\Sigma}h\colon
      \operatorname{comp}_{\Sigma}F
      &\longrightarrow \operatorname{comp}_{\Sigma}F'.
  \end{aligned}
  \tag{2.4.3.2}
\end{equation}

\VerdierSemanticAnchor{ch1-2-4-3-composition}
If $h$ and $h'$ are two composable morphisms of multiadditive functors, one
has:
\begin{equation*}
  \begin{aligned}
    \operatorname{comp}_{\Pi}hh'
      &=\operatorname{comp}_{\Pi}h\,\operatorname{comp}_{\Pi}h',\\
    \operatorname{comp}_{\Sigma}hh'
      &=\operatorname{comp}_{\Sigma}h\,\operatorname{comp}_{\Sigma}h'.
  \end{aligned}
\end{equation*}

\VerdierSemanticAnchor{ch1-2-4-3-sign-rule-ii}
The morphism $\operatorname{comp}_{\Pi}h$ (respectively,
$\operatorname{comp}_{\Sigma}h$) is a morphism of
$\mathbb Z[n]$-functors that obeys Sign Rule II (1.6.6).

\VerdierSemanticAnchor{ch1-2-4-4}
\noindent\textbf{2.4.4.}\quad
When the functor $F$ is contravariant in certain variables, i.e. when one
has equalities $C_i=(C_i')^\circ$, for $i\in J\subset[n]$, the functors
$\operatorname{comp}_{\Pi}F$ and $\operatorname{comp}_{\Sigma}F$ are,
modulo the identifications made in (2.4.1), contravariant in the
corresponding variables.

\VerdierSemanticAnchor{ch1-2-5}
\subsection{Homotopies}

\VerdierSemanticAnchor{ch1-2-5-1}
\noindent\textbf{2.5.1.}\quad
Let $X$ and $Y$ be two $n$-fold complexes of $C$, and let
$f,g\colon X\rightrightarrows Y$ be two morphisms of complexes. A
\emph{homotopy} connecting $f$ to $g$ is a family of morphisms

\VerdierSemanticAnchor{ch1-2-5-1-family}
\begin{equation*}
  a=\bigl(a^{\sigma,i}\colon X^{\sigma+e_i}\longrightarrow Y^\sigma,
    \ \sigma\in\mathbb Z[n],\ i\in[n]\bigr)
\end{equation*}
such that:

\VerdierSemanticAnchor{ch1-2-5-1-eq1}
\begin{equation}
  \begin{aligned}
    a^{\sigma+e_i-e_j,j}d_X^{\sigma,i}
      &=d_Y^{\sigma-e_j,i}a^{\sigma-e_j,j},
      &&\forall\sigma\in\mathbb Z[n],\
        \ \forall i,j\in[n],\ i\ne j,\\
    g^\sigma-f^\sigma
      &=\sum_{i\in[n]}\Bigl(
          a^{\sigma,i}d_X^{\sigma,i}
          +d_Y^{\sigma-e_i,i}a^{\sigma-e_i,i}\Bigr),
      &&\forall\sigma\in\mathbb Z[n].
  \end{aligned}
  \tag{2.5.1.1}
\end{equation}

\VerdierSemanticAnchor{ch1-2-5-1-notation}
We use the notation $a\colon f\rightsquigarrow g$ to denote homotopies $a$
connecting $f$ to $g$.

\VerdierSemanticAnchor{ch1-2-5-2}
\noindent\textbf{2.5.2.}\quad
Let $f,g,h\colon X\longrightarrow Y$ be three morphisms,
$a\colon f\rightsquigarrow g$ and $b\colon g\rightsquigarrow h$ two
homotopies. We denote by $b+a\colon f\rightsquigarrow h$ the homotopy

\VerdierSemanticAnchor{ch1-2-5-2-sum}
\begin{equation*}
  (b+a)^{\sigma,i}=b^{\sigma,i}+a^{\sigma,i},
  \qquad \sigma\in\mathbb Z[n],\ i\in[n].
\end{equation*}

\VerdierSemanticAnchor{ch1-2-5-3}
\noindent\textbf{2.5.3.}\quad
Let $f,g\colon X\rightrightarrows Y$ be two morphisms of complexes,
$a\colon f\rightsquigarrow g$ a homotopy, and let
$u\colon Y\longrightarrow Y'$ (respectively, $v\colon X'\longrightarrow X$)
be a morphism of complexes. We denote by
$u\star a\colon uf\rightsquigarrow ug$ (respectively,
$a\star v\colon fv\rightsquigarrow gv$) the homotopy

\VerdierSemanticAnchor{ch1-2-5-3-whiskering}
\begin{equation*}
  (u\star a)^{\sigma,i}=u^\sigma a^{\sigma,i}
  \qquad
  \text{(respectively, }(a\star v)^{\sigma,i}
    =a^{\sigma,i}v^{\sigma+e_i}\text{)}.
\end{equation*}
